\section{Introduction}
\label{sec:intro}

A skill is a reusable procedure, formula, lookup table or worked example supplied
in a language model's context to help it solve a task. Recent benchmarks show
that skills can improve performance \citep{su2026srabench, skillsbench2026},
but higher accuracy alone does not explain how a model uses them. Does a single
hidden-state vector carry the useful content, or does the model rely on states
distributed across the skill's token positions? How long does it continue to
read those states? We address these questions by intervening on the model's
representations and measuring how much of the skill's performance gain they
recover.

Our starting point is a contrast between behaviour and hidden-state geometry.
Across $1{,}100$ clinical-calculator instances, the correct skill improves
accuracy by $39.7$ percentage points over no skill, whereas a wrong skill of
the same token length changes it by $-2.9$ points ($95\%$ interval
$[-8.6,+2.9]$, paired calculator bootstrap). Yet in a synthetic task, correct and wrong skills produce
similar changes at the final prompt position: their state differences relative
to no skill have cosine similarity $0.91$--$0.95$ over layers $0$--$10$.
These observations motivate a causal question: does the shared state change
support the improvement, or is the content-specific difference responsible?

We subtract hidden states and inject each difference back into the model. On
the synthetic task, the change from no skill to a wrong skill recovers only
$0.098$ of the correct skill's gain, whereas the correct-minus-wrong difference
recovers $0.805$. Recovery also depends on answer format: the best state
replacement recovers $1.000$ for one-token multiple choice, $0.235$ for a short
number and $0.003$ for step-by-step reasoning. Success at the final prompt
position can therefore reflect transfer of a compact answer decision, while
failure need not imply that the skill lacks a transferable representation.

We find that \emph{states across the skill's own token span can transfer much
of its effect}. Replacing a wrong skill's states with the correct skill's at
one layer recovers $0.89$ of its gain on a screened subset of $135$ items from
$14$ clinical calculators. Halving the injection, permuting its positions or
substituting norm-matched noise recovers at most $0.05$. Truncating the matrix
of state differences tests how many directions support recovery. With the
positional mean preserved, $32$ leading directions recover $0.12$, $128$
recover $0.77$, and the bottom $64$ recover only $0.03$. Both position and the
retained directions matter for this intervention.

We compare models by the number of directions needed to reach half of the
measured recovery range. Reducing hidden width fourfold within Qwen changes
this value from $47$ to $43$, rather than the pre-registered proportional
prediction of approximately $12$; Mistral at the larger width gives $64$--$72$,
with intervals that overlap Qwen's.
Across layers, skill states stop transferring most of the effect while
attention to them still supports performance. This separation replicates
wherever both interventions are available. A few unusually large activations
explain an apparent drop in dimensionality, which therefore does not establish
compression. Our contributions are:
\begin{enumerate}\itemsep0pt\parskip0pt
\item \textbf{Locating a transferable skill representation.} Matched-control
subtraction isolates content-specific state differences, answer-format
comparisons reveal the limits of a single-position probe, and span injection
identifies a representation that recovers the skill's effect
(\S\ref{sec:decomp}--\ref{sec:whose}).
\item \textbf{Measuring the representation rank required for recovery.}
Truncating the injected matrix shows that recovery depends on multiple leading
directions. The measured rank curves depart from proportional width scaling
within Qwen (\S\ref{sec:rank}, \S\ref{sec:ladder}).
\item \textbf{Separating skill transfer from continued reading.} Span
replacement and blocking attention to skill tokens reveal distinct depth
boundaries. Interventions over multiple layers and checks for unusually large
activations help rule out two proposed explanations for the loss of transfer
(\S\ref{sec:depth}).
\end{enumerate}

